\section{Comparing insertions and their finite differences}
\label{sec:actual-insertion-derivative-relay}

The local equation comparison must be accompanied by two further
operations: representing a lower-grade insertion in a higher grade, and
representing a first finite difference in the next grade. We prove these
for a fixed simple gauge group on the same funded calibrated family.
We assume the quantitative service and positive reserve estimates of
the actual marked construction, together with its terminal activity
keepers and opened-scale response bounds, on this same history;
these are the hypotheses of Lemma~\ref{lem:actual-marked-antilles},
not additional conclusions of the present argument. For the difference
relay we also assume the derivative-shift purchase
$J_N\le\widehat\beta_N$ stated in its place below. The argument
uses the same terminal coefficient construction as the marked section,
but allows a local function whose coefficients vary with the cutoffs.

Let $0<a_N\le1$ be the lattice spacing and $s_N=\log(1/a_N)$. An insertion
$F_{N,p}(o)$ has engineering weight $e$ in the present calculation when
its unnormalized local function is
\[
                         f_{N,p,o}=a_N^eF_{N,p}(o).
\]
We require a fixed finite stencil containing zero, the regularity and
weighted left-derivative bounds of the actual top budget, and two
nonnegative subpower prices $c_{N,p},b_{N,p}$. The first bounds those
derivatives on the healthy region and includes the value bound there;
the second bounds $|f_{N,p,o}|$ on the whole configuration space. These
bounds are uniform in $o$. Locality, gauge invariance and the vanishing
Taylor conditions in the regularity hypothesis are imposed with the
same stencil and the same budget. We call these the protected local
bounds. Subpower has the ordered-cutoff meaning of
Section~\ref{sec:actual-equation-kernel-relay}.

\subsection{The physical section and a translated derivative stencil}

Write $\xi_{N,p}=X_{N,p}^{-1}x_{N,p}$ for the actual inverse terminal
coordinates of a label row. If this row is subpower in those coordinates,
its insertion has protected local bounds at weight $d$. Indeed its
unnormalized form is the physical section evaluated at
$S_{\theta_N^{-1}}\xi_{N,p}$, where the factor-degree scaling and
threshold are those used in the local section. The paid derivative bound
is
\[
             c_{N,p}=C_N|\xi_{N,p}|_1,
\]
with $C_N\ge1$ the top column allowance. For the global bound, compactness
of the fixed section's configuration cylinder gives a constant $A_d\ge1$
such that its value at any coefficient $v$ is at most $A_d|v|_1$.
To see that the constant is uniform in placement and period, expand in
the finitely many coefficient basis vectors on the fixed cover cylinder,
bound each continuous basis function on that compact cylinder, and take
their maximum. Every finite-volume placement is its pullback.
Thus one may use
\[
             b_{N,p}=A_d|S_{\theta_N^{-1}}\xi_{N,p}|_1.
\]
The column allowance and every fixed power of $\theta_N^{-1}$ are
subpower on the purchased schedule. Finite-dimensionality and the fixed
factor degrees therefore make both prices subpower. For a constant
reference label, the inverse terminal coordinates are subpower by
Section~\ref{ex-column-subpower}.

An insertion of weight $d$ with these bounds has them at any weight
$e\ge d$: multiply its unnormalized function and prices by
$a_N^{e-d}\le1$. This changes neither its stencil nor its locality.

For a fixed physical label $x$ of grade $d$, define its difference in
axis $\alpha$ by
\[
 D_{\alpha,N}O(x;o)=a_N^{-1}
             \{O_{N,p}(x;o+e_\alpha)-O_{N,p}(x;o)\}.
\]
At weight $d+1$ the apparently large factor $a_N^{-1}$ cancels exactly:
\begin{equation}\label{eq:relay-difference-weight}
 a_N^{d+1}D_{\alpha,N}O(x;o)
       =a_N^dO_{N,p}(x;o+e_\alpha)-a_N^dO_{N,p}(x;o).
\end{equation}
The common stencil is $(X+e_\alpha)\cup X$. Its derivative price must
also pay for moving the metric origin by one lattice unit.

Suppose the top budget uses derivative order $q_N\in\mathbb N$ and spatial power
$w_N\in\mathbb N$, so an allowed direction $A$ at origin $o$ satisfies
$\|A(z)\|\le\theta_N(1+\operatorname{dist}(z,o))^{w_N}$.
The triangle inequality gives
\[
 1+\operatorname{dist}(z,o)
        \le2\{1+\operatorname{dist}(z,o+e_\alpha)\}.
\]
Dividing all directions by $2^{w_N}$ makes them admissible at the
shifted origin. An ordered left derivative of length $q\le q_N$ is
multilinear in its directions, so restoring their size costs at most
\[
                         J_N=2^{w_Nq_N}.
\]
The protected prices for \eqref{eq:relay-difference-weight} are therefore
$(J_N+1)c_{N,p}$ and $2b_{N,p}$. We assume the selected-run
derivative-shift purchase $J_N\le\widehat\beta_N$ for all sufficiently
large $N$; it is a hypothesis on the selected delay, to be discharged by
the caller, and is not proved in this monograph. Given it, the hatted
coupling is subpower, so this extra price is subpower too. This order-and-power purchase is essential and
does not follow merely from the cancellation of $a_N^{-1}$.

\subsection{A terminal coefficient with a subpower total price}

We next encode any protected local function $f=a_N^eF_{N,p}(0)$ at
grade $e$. Enlarge its stencil by the stencil of the physical section
at grade $e$, and denote the resulting finite set by $D$. The readout
bound at this grade gives a fixed constant $K_e$ such that the leading
coefficient $z$ read from $f$ satisfies
\[
                              |z|_1\le q:=K_ec_{N,p}.
\]
Here and below the local coefficient includes the top threshold
normalization. Let $g$ be the physical local section for this same $z$,
and split $g=z_0+\xi$ into its scalar and nonconstant parts. Then
$y=f-g$ has zero readout. Its regular derivative price is at most
\[
                            h=c_{N,p}+C_Nq.
\]
Let $\mathsf N_N$ be the finite sum of the inverse threshold powers
over the grade-$e$ frame. The global section bound gives
$|g|,|\xi|\le A_e\mathsf N_Nq$. Thus the part left outside the healthy
masks has global bound
\[
                            r=b_{N,p}+2A_e\mathsf N_Nq.
\]

The exact decomposition is
\begin{equation}\label{eq:relay-protected-split}
 f=z_0+\mathbf1_{\rm healthy(X)}\xi
       +\mathbf1_{\rm healthy(D)}y
       +\bigl(y-\mathbf1_{\rm healthy(D)}y
                       +\xi-\mathbf1_{\rm healthy(X)}\xi\bigr),
\end{equation}
where $X\subseteq D$ is the physical section footprint. The bracketed
function vanishes when $D$ is healthy, is measurable and gauge invariant,
and depends on the same links and flags as the local construction. Its
absolute value is at most $r$, since each off-mask part is either zero
or its original value. The regular masked term has price $h$ and zero
readout. Apply the activity attachment construction used in the terminal
invitation to this bracketed function, using the actual top keeper.
The coefficient so obtained has leading value $z$ and exact density
\[
                         Z.\mathrm{density}=R_N f
                                            =a_N^eR_NF_{N,p}(0).
\]
No replacement of $R_N$ by its action without activities is made.

We give the total price explicitly. Write $\omega_N$ for the forced
allowance and $\nu_N>0$ for the support exponent, and put
\[
 A_N=e^{81\nu_N},\qquad M=2|D|.
\]
Eventually the placed stencil has covering number at most $81$ and
$|D|\omega_N\le1$; also $\sqrt{\omega_N}\le1$. The regular remainder
price from \eqref{eq:relay-protected-split} is bounded by
\[
       A_Nh+2|D|\sqrt{\omega_N}A_Nr
                              \le A_N(h+Mr)=:t_N.
\]
The forced attachment has reserved mass at most
\[
                         H_N=A_Ne^{162}r,
\]
because increasing the support exponent by two on covering number at
most $81$ costs at most $e^{162}$. Let $u_N>0$ be the selected attachment
reserve and $\zeta_{e,N}\ge0$ the grade-$e$ weight. Set
\begin{equation}\label{eq:relay-protected-total}
 Q_N=1+t_N+q,\qquad
 E_{N,p}=4(1+C_N)Q_N+H_N\zeta_{e,N}u_N^{-1}.
\end{equation}
Then $Q_N>0$, the leading coefficient plus regular remainder costs at
most $Q_N$, the first summand pays its required terminal cone allowance,
and $H_N\zeta_{e,N}\le u_NE_{N,p}$ pays the forced attachment. The exact
profile identity comes from the same attachment; it is preserved in
the coefficient density just displayed.

Every factor in \eqref{eq:relay-protected-total} is subpower. For $q,h,r$
this follows from the protected prices, the column bound and the finite
threshold sum. The support-exponent purchase makes $A_N$ subpower.
The selected reserve record gives subpower bounds for $\zeta_{e,N}$
and $u_N^{-1}$. Finite sums and products therefore make $E_{N,p}$
subpower. Continuity of the original insertion completes the terminal
input required by the descent; continuity across a mask boundary is
not asserted for the individual masked pieces.

\subsection{Descent and the comparison labels}

Apply the same-history backward coefficient construction of
Lemma~\ref{lem:actual-marked-antilles} to this terminal coefficient.
Retain also the served-representation and terminal-attachment inputs
of Propositions~\ref{prop:actual-marked-figment}
and~\ref{prop:terminal-marked-voice} on this history.
Its bottom primary determines a reference label $v_{N,p}$ of grade
$e$. Item (4) of that lemma, before taking the last-step purse,
gives
\[
 32(1+C_N)^2|Z.x-X_{N,p}^{-1}v_{N,p}|_1
                    \le2(1+C_N)Q_N\le E_{N,p}.
\]
Also $|Z.x|_1\le Q_N\le E_{N,p}$. Since $32(1+C_N)^2\ge1$, the
triangle inequality yields
\[
                         |X_{N,p}^{-1}v_{N,p}|_1\le2E_{N,p}.
\]
The comparison label is thus subpower in its actual inverse terminal
coordinates. Choose these labels at every cutoff where the eventual
terminal construction applies, and arbitrary labels elsewhere.

The same descent and finite-census argument give, for
$\operatorname{ref}(e)\le k<N$, $m\le N$, and every opened pair $(m,k)$,
\[
 \mathcal R_{m,N,p}\{F_{N,p}(0)-O_{N,p}(v_{N,p};0)\}
                     \le c_{e,m,k}E_{N,p}e^{-s_N/8}.
\]
Replacing the difference by its negative preserves its root norm.
The fixed-stencil and common-tiling argument of
Section~\ref{sec:actual-equation-kernel-relay} converts this into a
vanishing discrepancy. It is uniform in the required translated
placements in the ordered sense. The finite-cutoff density identity,
the descent, and this discrepancy all use one original history.

Apply this result first to the fixed insertion of grade $d$, at any
weight $e\ge d$, and then to its difference in axis $\alpha$, at weight
$d+1$. We obtain in each case comparison labels of the requested grade,
subpower inverse terminal coordinates, and a discrepancy tending to
zero against every fixed separated spectator test. These are the
inclusion and derivative comparison inputs used to build the physical
filtration and field hierarchy. Together with the equation comparison
of the preceding section, they supply the three local comparison
clauses on every subsequent refinement of the ordered extraction:
eventual inequalities, subpower bounds and limits all persist on such
refinements.
